\scanpage{83}
% archive sheet number: 82
% block text | top addition and lines 1-3
\inserted{Supposons que $M$ soit un groupe abélien.} La façon habituelle de construire des éléments de $M_f$ (du moins en supposant $\Gamma$ \uline{fini}) est \deleted{\unsure{par}} en prenant des moyennes
% block display | line 4
\[
N_f m=\sum_{\unsure{\gamma|\sigma}\in\Gamma}\gamma_f m=\sum_{\gamma\in\Gamma}f(\gamma)\gamma\cdot m
\]
% block text | lines 5-8
Supposons en particulier que $G$ soit l'ens des éléments inversibles d'un \uline{anneau} $A$, et que $M$ soit un $A$-module. Soit \deleted{\unsure{$N_f$ l'application de $A$|$N_f$ l'opérateur de $A$}}\inserted{pour \deleted{\unsure{$\forall$}} $u\in A$} \deleted{défini par}
% block display | line 8
\[
N_f u=\sum_{\gamma\in\Gamma}f(\gamma)\gamma(u)
\]
% block margin note | right of line 8
\margin{$\boxed{N_f(u)\deleted{\unsure{\in}}\in A_f^\Gamma}$}
% block text | line 9
Alors
% block display | line 9
\[
\deleted{\unsure{N_f u|N_f m}}xM^\Gamma\unsure{\subset|\in}M_f^\Gamma
\]
% block margin note | right of lines 9-10
\margin{En particulier, \unsure{$=$}}
% block display | line 10
\[
\text{i.e.}\qquad f(\gamma)\gamma(x)=\deleted{\unsure{f(\gamma)}}x\qquad\text{\unsure{où|si} }x=N_f u
\]
% block text | lines 11-14
[Si le groupe $\Gamma$ est infini, on essaie de faire des moyennes infinies, quand on arrive à leur donner un sens. C'est la \deleted{\unsure{méthode}} construction des fonctions \unsure{zeta-fuchsiennes|zêta-fuchsiennes|thêta-fuchsiennes} de Poincaré.]

% block text | small addition between lines 14 and 15
\inserted{Si on peut s'arranger pour trouver $u$ tel que $x$ soit inversible, on aura $f(\gamma)=x\gamma(x)^{-1}$ donc $f\sim0$}

% block text | line 15
\uline{Exemple} Supposons que $\Gamma$ \uline{\unsure{fini}} opère \uuline{fidèlement} dans \inserted{un} \uline{corps} $K$.

% block text | lines 16-17
Soit $A$ une algèbre de dim finie sur $K$, et supposons que $\Gamma$ opère dans $A$, de façon que
% block display | line 18
\[
\gamma(\lambda a)=(\gamma\cdot\lambda)(\gamma\cdot a)
\]
% block text | lines 19-23
Je dis qu'alors $H^1(\Gamma,G)\simeq(e)$, où $G$ est le groupe des éléments inversibles de $A$. D'après ce qui précède, il suffit de trouver un $\unsure{u|a}\in A$ tel que $\sum f(\gamma)\cdot\gamma u$ soit inversible, ce qui s'exprime par une \unsure{équation}.
